\documentclass[14pt]{extreport}

% Основные пакеты
\usepackage[english, russian]{babel}    % Поддержка русского и английского языков
\usepackage{graphicx}                   % Для вставки изображений
\usepackage{import}                     % Для продвинутого импорта (например, pdf_tex из Inkscape)
\usepackage{subcaption}                 % Для создания подписей к фигурам
\usepackage{float}                      % Для улучшенного позиционирования фигур
\usepackage{geometry}                   % Для управления отступами на странице
\usepackage{amsmath}                    % Для расширенной поддержки математики
\usepackage{amsfonts}                   % Математические шрифты
\usepackage{amssymb}                    % Математические символы
\usepackage{tabu}                       % Улучшенные таблицы
\usepackage{booktabs}                   % Профессиональные таблицы с вертикальными разделителями
\usepackage[dvipsnames]{xcolor}         % Расширенная поддержка цветов
\usepackage{longtable}                  % Для таблиц, занимающих несколько страниц
\usepackage{makecell}                   % Для улучшенного форматирования ячеек таблиц
\usepackage{multirow}                   % Для объединения строк в таблицах
\usepackage{tabularray}                 % Для современных таблиц
\usepackage{varwidth}                   % Для переменной ширины блоков
\usepackage{hyperref}                   % Для гиперссылок
\usepackage{verbatim}                   % Для вставки кода
\usepackage{listings}                   % Для оформления листингов кода
\usepackage{amsmath}
\usepackage{listings}
\newcommand{\outputbox}[1]{
	\fboxsep=5pt
	\colorbox{lightgray!20}{\texttt{#1}}
}
\usepackage{listings}
\usepackage[T2A]{fontenc} % Для поддержки кириллицы
\usepackage[utf8]{inputenc}
\usepackage[russian]{babel}
\usepackage{tcolorbox}
\tcbuselibrary{listings}
% Настройка графики
\graphicspath{ {./images/} }            % Путь к папке с изображениями

% Настройка отступов на странице
\geometry{left=2cm, right=2cm, bottom=2cm, top=2cm}

% Настройка гиперссылок
\hypersetup{
	linkcolor=blue,                     % Цвет ссылок
	filecolor=magenta,                  % Цвет ссылок на файлы
	urlcolor=cyan                       % Цвет URL-адресов
}

\definecolor{codegreen}{rgb}{0,0.6,0}  % Цвет для комментариев
\definecolor{codegray}{rgb}{0.5,0.5,0.5}  % Цвет для номеров строк
\definecolor{codepurple}{rgb}{0.58,0,0.82}  % Цвет для строк
\definecolor{backcolour}{rgb}{0.95,0.95,0.92}  % Цвет фона

% Определим стиль для выделения переменных
\lstdefinestyle{codestyle}{
	backgroundcolor=\color{backcolour},  % Цвет фона
	commentstyle=\color{codegreen}\itshape,  % Стиль для комментариев
	keywordstyle=\color{magenta}\bfseries,  % Стиль для ключевых слов
	numberstyle=\tiny\color{codegray},  % Стиль для номеров строк
	stringstyle=\color{codepurple},  % Стиль для строк
	identifierstyle=\color{black},  % Стиль для переменных (выделим их синим)
	basicstyle=\ttfamily\footnotesize,  % Размер шрифта для кода
	breakatwhitespace=false,  % Не разрывать строки в пробелах
	breaklines=true,  % Разбивать длинные строки
	captionpos=b,  % Позиция заголовка (внизу)
	keepspaces=true,  % Сохранять пробелы
	numbers=left,  % Позиция номеров строк
	numbersep=10pt,  % Расстояние между кодом и номерами строк
	showspaces=false,  % Не показывать пробелы
	showstringspaces=false,  % Не показывать пробелы в строках
	showtabs=false,  % Не показывать табуляции
	tabsize=4,  % Размер табуляции
	frame=single,  % Добавить рамку вокруг кода
	rulecolor=\color{black},  % Цвет рамки
	backgroundcolor=\color{backcolour},  % Цвет фона рамки
	xleftmargin=15pt,  % Отступ слева
	xrightmargin=15pt,  % Отступ справа
}


\lstset{style=codestyle}
\lstset{extendedchars=\true}
\graphicspath{{../images/}}


\begin{document}
	
	\begin{titlepage}
		\centering
		\vspace*{1cm}
		
		{\large Министерство науки и высшего образования Российской Федерации}\\
		{\large ФЕДЕРАЛЬНОЕ ГОСУДАРСТВЕННОЕ АВТОНОМНОЕ ОБРАЗОВАТЕЛЬНОЕ УЧРЕЖДЕНИЕ ВЫСШЕГО ОБРАЗОВАНИЯ «НАЦИОНАЛЬНЫЙ ИССЛЕДОВАТЕЛЬСКИЙ УНИВЕРСИТЕТ ИТМО»}\\
		{\large (УНИВЕРСИТЕТ ИТМО)}\\
		
		\vspace{1cm}
		
		\textbf{{\Huge ОТЧЕТ}\\
			{\Huge ПО ЛАБОРАТОРНОЙ РАБОТЕ №4}}\\
		
		\vspace{1cm}
		
		{\LARGE По дисциплине\\
			<<Математическая Статистика>>}\\
		Вариант 2
		
		\vspace{2cm}
		
		{\Large Студенты:}\\
		Охрименко Eва\\
		Даниил Буцкий
		
		\vspace{2cm}
		
		{\Large Проверил:}\\
		Шкваренко Андрей Алексеевич\\
		
		\vspace{2cm}
		
		{\large г. Санкт-Петербург}\\
		{\large 2025}
	\end{titlepage}
	
	\tableofcontents
	\newpage
	
	\setcounter{section}{0} % Сбрасываем счетчик разделов
	
	\section{Task 1}
	
	\section{Task 2}
	
	В этом задании нужно определить, влияет ли национальность студентов на их общий балл за экзамены. Другими словами — есть ли группы, которые сдают экзамены статистически значимо лучше или хуже других. \\
	
	
	Для начала создадим несколько массивов инициализированных по этническому признаку. И засунем туда суммарные баллы студентов именно этой национальности. Вот код, реализующий это:
	
	\begin{lstlisting}[language=python, caption={Инициализация массивов с баллами}]
f = open('exams_dataset.csv')

countA = []
countB = []
countC = []
countD = []
countE = [] 
continueCount = 0

for line in f:

	continueCount += 1
	if continueCount == 1:
		continue
	
	items = line.strip().replace('"', '').split(',')
	
	etnic = items[1]
	points = sum(map(int, items[-3:])) 
	
	if etnic[-1] == 'A':
		countA.append(points)
	elif etnic[-1] == 'B':
		countB.append(points)
	elif etnic[-1] == 'C':
		countC.append(points)
	elif etnic[-1] == 'D':
		countD.append(points)
	elif etnic[-1] == 'E':
		countE.append(points)
\end{lstlisting}
	
	
	Для того, чтобы сделать модель однофакторного дисперсионного анализа нам надо найти:
	
	\begin{itemize}
		\item $\mu$ --- общее среднее по всем группам
		\item \texttt{SSB} --- cумма квадратов между группами
		\item \texttt{SSW} --- cумма квадратов внутри групп
		\item степени свободы
		\item \texttt{F}-статистику
		\item критическое значение $\mathtt{F\_crit}$, чтобы сравнить с \texttt{F}
	\end{itemize}
	
	Итак, начнем. Для начала расчитаем общее среднее по всем группам:
	\[
	\mu = \frac{\sum_{i=1}^{k} \sum_{j=1}^{n_i} Y_{ij}}{N}
	\]
	
	\begin{lstlisting}[language=python, caption=mean по всем группам]
allScores = countA + countB + countC + countD + countE
mu = sum(allScores) / len(allScores)
\end{lstlisting}
	
	
\noindent
\fbox{
	\begin{minipage}{\dimexpr\linewidth-2\fboxsep-2\fboxrule\relax}
		\vspace{0.5em} 
		\texttt{mu = 202.404}
		\vspace{0.5em} 
	\end{minipage}
}
\\

	Теперь найдем \texttt{SSB}:
	\[
	\mathtt{SSB} = \sum_{i=1}^{k} n_i (\bar{Y}_i - \mu)^2
	\]
	
	
	\begin{lstlisting}[language=python, caption=\texttt{SSB}]
groupMeans = [
	sum(countA)/len(countA), 
	sum(countB)/len(countB), 
	sum(countC)/len(countC), 
	sum(countD)/len(countD), 
	sum(countE)/len(countE)  
]

groupSizes = [len(countA), len(countB), len(countC), len(countD), len(countE)]

SSB = sum(
	n * (groupMean - mu) ** 2 
	for n, groupMean in zip(groupSizes, groupMeans)
)
\end{lstlisting}

\noindent
\fbox{
	\begin{minipage}{\dimexpr\linewidth-2\fboxsep-2\fboxrule\relax}
		\vspace{0.5em} 
		\texttt{SSB	 $\approx$ 112381.64}
		\vspace{0.5em} 
	\end{minipage}
}
\\

	
	Расчитаем \texttt{SSW}:
	\[
	\texttt{SSW} = \sum_{i=1}^{k} \sum_{j=1}^{n_i} (Y_{ij} - \bar{Y}_i)^2
	\]
	
	\begin{lstlisting}[language=python, caption=\texttt{SSW}]
groups = [countA, countB, countC, countD, countE]
SSW = 0.0
for group, mean in zip(groups, groupMeans):
	for score in group:
		SSW += (score - mean) ** 2
\end{lstlisting}	
	
\noindent
\fbox{
	\begin{minipage}{\dimexpr\linewidth-2\fboxsep-2\fboxrule\relax}
		\vspace{0.5em} 
		\texttt{SSW $\approx$ 1953413.14}
		\vspace{0.5em} 
	\end{minipage}
}
\\
	
	
	
	Теперь посчитаем степени свободы:
	
	\begin{align*}
		df_{total} &= N - 1 \\
		df_{between} &= k - 1 \\
		df_{within} &= N - k
	\end{align*}
	
	Где:
	\begin{itemize}
		\item $N$ - общее количество наблюдений
		\item $k$ - количество групп
		\item $df_{total}$ - общие степени свободы
		\item $df_{between}$ - степени свободы между группами
		\item $df_{within}$ - степени свободы внутри групп
	\end{itemize}
	
\begin{lstlisting}[language=python, caption=Степени свободы]
N = sum(len(group) for group in groups)  
k = len(groups)                           

dfTotal = N - 1
dfBetween = k - 1
dfWithin = N - k
\end{lstlisting}

\noindent
\fbox{
	\begin{minipage}{\dimexpr\linewidth-2\fboxsep-2\fboxrule\relax}
		\vspace{0.5em} 
		\texttt{dfTotal = 999 \\
			dfBetween = 4 \\
			dfWithin = 995
		}
		\vspace{0.5em} 
	\end{minipage}
}
\\


Далее посчитаем \texttt{F}-статистику
\[
\texttt{F} = \frac{MSB}{MSW},
\]

Где: 
\begin{itemize}
	\item $MSB = \frac{SSB}{df_{between}}$ 
	\item $MSW = \frac{SSW}{df_{within}}$ 
\end{itemize}


	\begin{lstlisting}[language=python, caption=\texttt{F}-статистика]
MSB = SSB / dfBetween 
MSW = SSW / dfWithin
F = MSB / MSW 
\end{lstlisting}

\noindent
\fbox{
	\begin{minipage}{\dimexpr\linewidth-2\fboxsep-2\fboxrule\relax}
		\vspace{0.5em} 
		\texttt{F $\approx$ 14.31}
		\vspace{0.5em} 
	\end{minipage}
}
\\

Теперь найдем $\mathtt{F\_crit}$ кодом. Возьмем $\alpha = 0.05$ - стандартный уровень значимости.

	\begin{lstlisting}[language=python, caption=\texttt{F}-статистика]
alpha = 0.05
F_crit = f.ppf(1 - alpha, dfBetween, dfWithin)
\end{lstlisting}

\noindent
\fbox{
	\begin{minipage}{\dimexpr\linewidth-2\fboxsep-2\fboxrule\relax}
		\vspace{0.5em} 
		\texttt{F\_crit $\approx$ 2.38}
		\vspace{0.5em} 
	\end{minipage}
}
\\



\textbf{Итог:} \\
В результате однофакторного дисперсионного анализа получилось

\begin{center}
	\texttt{F} > $\mathtt{F\_crit}$.
\end{center}
Это означает, что количество баллов зависит от национальности.	

\end{document}